\subsection{Nagata 环}

定义 2. --- 称环 A 为 Nagata 环，如果它是 Noether 环，并且对于 A 的每个素理想 p，Noether 整环 A/p 都是日本环（第 1 目，定义 1）。

例. --- 1) 域上的每个有限生成代数都是 Nagata 环（第 1 目，例）。

\begin{enumerate}
\def\labelenumi{\arabic{enumi})}
\setcounter{enumi}{1}
\item
  每个完备局部 Noether 环都是 Nagata 环（第 2 目，定理 2）。
\item
  环 Z 是 Nagata 环（第 1 目，例以及命题 1 的推论）。
\item
  可以证明（习题 30），Nagata 环上的每个有限生成代数都是 Nagata 环。
\end{enumerate}

命题 4. --- 设 A 是一个 Nagata 环。

\begin{enumerate}
\def\labelenumi{\alph{enumi})}
\item
  每个有限 A-代数都是 Nagata 环。
\item
  对于 A 的每个乘性子集 S，环 \(\mathbf{S}^{-1}\mathbf{A}\) 都是 Nagata 环。
\item
  设 B 是一个有限 A-代数，\(\rho: A \to B\) 为典范同态。对于 B 的每个素理想 p，环 B/p 是日本环 \(A/\rho^{-1}(p)\) 上的一个有限代数，故为日本环（第 1 目，命题 2）。
\item
  设 q 是 \(S^{-1}A\) 的一个素理想；则存在 A 的素理想 p 使得 \(q = S^{-1}p\) 。环 ( \(S^{-1}A\) )/q 是日本环 A/p 的一个分式环，故为日本环（第 1 目，注记 2）。
\end{enumerate}

定理 3 (Zariski-Nagata). --- 设 A 是一个半局部 Noether 环。下列条件等价：

\begin{enumerate}
\def\labelenumi{(\roman{enumi})}
\item
  A 是 Nagata 环；
\item
  对于 A 的每个素理想 \(\mathfrak{p}\)，\(\kappa(\mathfrak{p})\) -代数 \(\kappa(\mathfrak{p}) \otimes_{\mathbf{A}} \hat{\mathbf{A}}\) 是可分的；
\item
  对于每个约化 A-代数 R，环 \(R \otimes_{A} \hat{A}\) 都是约化的。
\end{enumerate}

我们先证明条件 (ii) 与 (iii) 的等价性。蕴涵 (iii) \(\Rightarrow\) (ii) 是显然的；反之，设 A 满足条件 (ii)。于是，对于每个是域的 A-代数 K，环 K \(\otimes_{\mathbf{A}}\hat{\mathbf{A}}\) 都是约化的。现设 C 是一个约化的有限生成 A-代数；环 C 作为 Noether 环，同构于域的有限乘积 K \(_1\) \(\times\) \(\cdots\) \(\times\) K \(_n\) 的一个子环（IV，§ 2，第 5 目，命题 10）；由于 \(\hat{\mathbf{A}}\) 在 A 上平坦，环 C \(\otimes_{\mathbf{A}}\hat{\mathbf{A}}\) 同构于约化环 \(\prod_{i}(\mathrm{K}_{i}\otimes_{\mathbf{A}}\hat{\mathbf{A}})\) 的一个子环，故为约化的。最后，设 R 是任一约化 A-代数；则 R 是其有限生成子代数的过滤族 \((\mathrm{C}_{\alpha})\) 的并，且 \(R \otimes_{A} \hat{A}\) 是约化环的过滤族 \((\mathrm{C}_{\alpha} \otimes_{A} \hat{A})\) 的归纳极限，故为约化的。

我们来证明 (ii) 蕴涵 (i)。设 p 是 A 的一个素理想；环 A/p 的分式域 K 等同于 \(\kappa(p)\) ，而 K-代数 \(K \otimes_{A/p} (\widehat{A/p})\) 等同于 \(\kappa(p) \otimes_{A/p} \widehat{A}/p\widehat{A}\) ，从而等同于 \(\kappa(p) \otimes_{A} \widehat{A}\) 。若 \(\kappa(p) \otimes_{A} \widehat{A}\) 是一个 \(\kappa(p)\) -可分代数，则环 A/p 是日本环（第 2 目，命题 3）。

我们对 dim(A) 作归纳来证明蕴涵 (i) \(\Rightarrow\) (ii)。当 dim(A) = 0 时这是显然的，因为此时 A 是 Artin 环，从而是完备的。设 \(n\) 是一个整数 \(>0\) ；考虑如下的假设：

\((\mathbf{R}_n)\left\{
\begin{array}{ll}\text{对于每个局部诺特 Nagata 环 C，其维数满足 } < n\text{，且每个素}\\ \text{理想 } \mathfrak{r}\text{均为 C 的理想，则该环 }\kappa(\mathfrak{r})\otimes_{\mathrm{C}}\hat{\mathrm{C}}\text{是约化的。} \end{array}
\right.\)

设 A 是一个维数为 n 的半局部 Noether Nagata 环，p 是 A 的一个素理想，L 是域 \(\kappa(\mathfrak{p})\) 的一个有限扩张；在假设 (R \(_{n}\) ) 之下，只需证明环 \(L \otimes_{A} \hat{A}\) 是约化的。用 B 表示 A/p 在 L 中的整闭包；由于 A/p 是日本环，B 是一个有限 A-代数，因而是半局部 Nagata 环（命题 4）。设 m \(_{1}\) , \ldots, m \(_{r}\) 是 B 的极大理想；环 \(L \otimes_{A} \hat{A}\) 等同于 \(B \otimes_{A} \hat{A}\) 的一个分式环，而后者等同于诸局部环 B \(_{m_{i}}\) 的完备化之乘积（第 3 目，引理 1）。因此只需证明，对于 B 的每个极大理想 m，环 B \(_{m}\) 都是约化的（II，§ 2，第 6 目，命题 17）。环 B \(_{m}\) 是局部的、整闭的且是 Nagata 的（命题 4），并且有 dim(B \(_{m}\) ) ≤ dim(B) ≤ dim(A) = n（VIII，§ 1，第 3 目，命题 6 以及 § 2，第 3 目，定理 1）。改变记号，我们归结为在假设 (R \(_{n}\) ) 之下证明：对于每个维数 ≤ n 的局部 Noether 整闭 Nagata 环 A，环 \(\hat{A}\) 是约化的，也就是说（第 3 目，引理 2），\(\hat{A}_{p'}\) 对每个素理想 p′ ∈ Ass(\(\hat{A}\)) 都是约化的。由于当 dim(A) = 0 时这是直接的，我们可以假设 dim(A) > 0。于是设 x 是 m \(_{A}\) 的一个非零元素，q' 是 \(\hat{A}\) 的一个素理想，它在包含 \(x\hat{A} + p'\) 的素理想中极小；由于 \(\hat{A}_{p'}\) 等同于环 \(\hat{A}_{q'}\) 的一个分式环，只需证明后者是约化的（II，§ 2，第 6 目，命题 17）。由引理 3，理想 q' 与 \(\hat{A}\)-模 \(\hat{A}/x\hat{A}\) 相伴；由于 \(\hat{A}\) 在 A 上平坦，q' 在 A 中的逆像 q 与 A-模 A/xA 相伴（IV，§ 2，第 6 目，定理 2 的推论 1）。由于环 A 被假设为整闭的，这蕴涵 q 的高度为 1（VII，§ 1，第 6 目，命题 10），从而环 \(A_{q}\) 是一个离散赋值环（同前，第 3 目，定理 2 的推论以及第 6 目，定理 4）。由于 A/q 是一个维数 < n 的 Nagata 环，环 \(\kappa(q) \otimes_{A/q} \hat{A}/q\hat{A}\) 由假设 (R \(_{n}\) ) 是约化的。与它同构的环 \(\kappa(q) \otimes_{A} \hat{A}\) 是约化的，因而环 \(\kappa(q) \otimes_{A_{q}} \hat{A}_{q'}\) 作为它的一个分式环也是约化的。于是可以把第 3 节的引理 4 应用于典范同态 \(A_{q} \to \hat{A}_{q'}\)，从而得出环 \(\hat{A}_{q'}\) 是约化的，这就是我们想要证明的。定理 3 至此证毕。

推论 1. --- 局部且约化的 Nagata 环的完备化是约化的。

事实上，只需在定理 3 的 (iii) 中取 R = A。

推论 2 (Chevalley). --- 设 A 是域上的一个有限型约化代数，p 是 A 的一个素理想。局部环 A \(_{p}\) 的完备化是约化的。

由于 A 是约化的，局部环 \(A_{p}\) 是约化的；此外，A 是一个 Nagata 环（例 1），因而 \(A_{p}\) 是 Nagata 环（命题 4），而把推论 1 应用于环 \(A_{p}\) 即得推论 2。

推论 3. --- 设 k 是一个特征为 0 的域，A 是一个局部 Noether k-代数。A 为 Nagata 环的必要充分条件是：对于 A 的每个素理想 p，环 (A/p) 都是约化的。

事实上，由于诸域 \(\kappa(p)\) 的特征为 0，说代数 \(\kappa(p) \otimes_{A} \hat{A} = \kappa(p) \otimes_{A/p} (\widehat{A/p})\) 是约化的，与说它们是可分的是等价的（A，V，第 117 页，定理 1），这表明所述条件是充分的（定理 3，(ii) \(\Rightarrow\) (i)）；而且它也是必要的（定理 3，(i) \(\Rightarrow\) (iii)，取 \(R = A/p\) ）。

\section*{附录}
\phantomsection
\addcontentsline{toc}{section}{附录}

\refstepcounter{appendixitem}
\subsection*{\theappendixitem.\ 局部环的归纳极限}
\phantomsection
\addcontentsline{toc}{subsection}{\theappendixitem.\ 局部环的归纳极限}

设 I 是一个非空的右过滤预序集，\((\mathbf{A}_{\alpha}, \varphi_{\beta\alpha})\) 是关于 I 的一个环的归纳系统。假设对每个 \(\alpha \in I\) ，环 \(A_{\alpha}\) 都是局部的，其极大理想为 \(m_{\alpha}\) ，诸同态 \(\varphi_{\beta\alpha}\) 都是局部且平坦的，并且有 \(\varphi_{\beta\alpha}(m_{\alpha}) A_{\beta} = m_{\beta}\) 对 \(\beta \geqslant \alpha\) 成立。用 A 表示诸 \(A_{\alpha}\) 的归纳极限，且对每个 \(\alpha \in I\) ，设 \(\varphi_{\alpha}: A_{\alpha} \to A\) 为典范同态。

命题 1. --- a) 环 A 是局部的，其极大理想为 \(\mathfrak{m} = \varinjlim \mathfrak{m}_{\alpha}\) 。对每个 \(\alpha \in \mathbf{I}\) ，同态 \(\varphi_{\alpha}\) 是局部且平坦的，且有 \(\varphi_{\alpha}(\mathfrak{m}_{\alpha})\) A = \(\mathfrak{m}\) 。

\begin{enumerate}
\def\labelenumi{\alph{enumi})}
\setcounter{enumi}{1}
\item
  若 \(\mathbf{A}_{\alpha}\) 对每个 \(\alpha \in \mathbf{I}\) 都是 Noether 的，则 \(\mathbf{A}\) 是 Noether 的。
\item
  令 \(m = \varinjlim m_{\alpha}\) ；它是 A 的一个理想。商环 A/m 是诸域 \(A_{\alpha}/m_{\alpha}\) 的归纳极限，因此是一个域（A，I，第 116 页，命题 3）。此外，A - m 的每个元素在 A 中可逆：事实上，设 \(x \in A - m\) ；存在 \(\alpha \in I\) 和 \(\xi \in A_{\alpha}\) 使得 \(x = \varphi_{\alpha}(\xi)\) ；我们有 \(\xi \notin m_{\alpha}\) ，因而 \(\xi\) 在 \(A_{\alpha}\) 中可逆，而 x 在 A 中可逆。于是，A 是一个以 m 为极大理想的局部环。设 \(\alpha \in I\) 。由关系式 \(\varphi_{\beta\alpha}(m_{\alpha}) A_{\beta} = m_{\beta}\) （\(\beta \geqslant \alpha\) ）取归纳极限即得 \(\varphi_{\alpha}(m_{\alpha}) A = m\) ；最后，由 I，§ 2，第 7 目，命题 9，同态 \(\varphi_{\alpha}\) 是平坦的。
\item
  设 \(\hat{\mathbf{A}}\) 是 \(\mathbf{A}\) 关于 m-进拓扑的分离完备化，\(\pi\) 是从 \(\mathbf{A}\) 到 \(\hat{\mathbf{A}}\) 的典范映射。假设诸环 \(\mathbf{A}_{\alpha}\) 都是 Noether 的。固定 \(\alpha \in \mathbf{I}\) ，证明环 \(\hat{\mathbf{A}}\) 在 \(\mathbf{A}_{\alpha}\) 上是 Noether 且平坦的。由假设，\(m_{\alpha}\) 是 \(\mathbf{A}_{\alpha}\) 的一个有限生成理想，故 \(m = \varphi_{\alpha}(m_{\alpha})\) 。\(m\) 是 \(\mathbf{A}\) 的一个有限生成理想。由此得出 \(\hat{m}\) （\(\hat{\mathbf{A}}\) 的极大理想）等于 \(m\hat{\mathbf{A}}\) （III，§ 2，第 12 目，命题 16 的推论 2 以及第 13 目，命题 19），从而是有限生成的。于是，环 \(\hat{\mathbf{A}}\) 是 Noether 的（同前，第 10 目，定理 2 的推论 5）。另一方面，对每个 \(n \in \mathbf{N}\) ，商 \(\hat{\mathbf{A}}/\hat{m}^{n}\) 同构于 \(\mathbf{A}/m^{n}\) （同前，第 12 目，命题 16 的推论 2 以及公式 (21)），这就意味着 \(\hat{\mathbf{A}}/\pi \circ \varphi_{\alpha}(m_{\alpha}^{n})\) \(\hat{\mathbf{A}}\) 同构于 \(\mathbf{A} \otimes_{\mathbf{A}_{\alpha}} (\mathbf{A}_{\alpha}/m_{\alpha}^{n})\) ；由于 \(\mathbf{A}\) 是平坦 \(\mathbf{A}_{\alpha}\)-模，\((\mathbf{A}_{\alpha}/m_{\alpha}^{n})\) -模 \(\hat{\mathbf{A}}/\pi \circ \varphi_{\alpha}(m_{\alpha}^{n})\) \(\hat{\mathbf{A}}\) 对每个 \(n \in \mathbf{N}\) 都是平坦的。由 III，§ 5，第 4 目，命题 2，\(\mathbf{A}_{\alpha}\) -模 \(\hat{\mathbf{A}}\) 关于 \(m_{\alpha}\) 是理想分离的；由同前，第 2 目，定理 1，\(\mathbf{A}_{\alpha}\) -模 \(\hat{\mathbf{A}}\) 因而是平坦的。取归纳极限即得 \(\hat{\mathbf{A}}\) 在 \(\mathbf{A}\) 上是（忠实）平坦的（I，§ 2，第 7 目，命题 9），从而 \(\mathbf{A}\) 是 Noether 的（I，§ 3，第 5 目，命题 8 的推论）。
\end{enumerate}

\refstepcounter{appendixitem}
\subsection*{\theappendixitem.\ 局部环的膨胀}
\phantomsection
\addcontentsline{toc}{subsection}{\theappendixitem.\ 局部环的膨胀}

设 A 是一个局部环。

我们用 A{]}X{[} 表示多项式环 A{[}X{]} 在素理想 \(m_{A}A[X]\) 处的局部环。这是一个以 \(m_{A}A{]}X{[}\) 为极大理想的局部环；典范同态 A → A{]}X{[} 是局部且平坦的，且 A{]}X{[} 的剩余域是由 X 的类生成的 \(\kappa_{A}\) 的纯扩张。

引理 1. — 设 \(\mathbf{P} \in \mathbf{A}[\mathbf{X}]\) 是一个首一多项式，其像 \(\overline{\mathbf{P}}\) 在 \(\kappa_{\mathbf{A}}[\mathbf{X}]\) 中不可约。则 A-代数 \(\mathbf{B} = \mathbf{A}[\mathbf{X}] / (\mathbf{P})\) 是局部的且在 \(\mathbf{A}\) 上有限，其极大理想为 \(\mathfrak{m}_{\mathbf{A}}\mathbf{B}\)；典范同态 \(\rho: \mathbf{A} \to \mathbf{B}\) 是局部且平坦的，剩余域扩张 \(\kappa_{\mathbf{A}} \to \kappa_{\mathbf{B}}\) 是代数的且由 \(x\) （\(\mathbf{X}\) 的类）所生成，而 \(x\) 在 \(\kappa_{\mathbf{A}}\) 上的最小多项式是 \(\overline{\mathbf{P}}\) 。

由于多项式 P 是首一的，A-模 B 是有限型的自由模（A，IV，第 10 页）。环 B/m \(_{A}\) B 等同于 \(\kappa_{A}[X]/(\overline{P})\) ，因而是一个域；于是理想 m \(_{A}\) B 是极大的。设 q 是 B 的一个极大理想；则理想 \(\rho^{-1}(q)\) 是极大的（V，§ 2，第 1 目，命题 1）；因此我们有 \(\rho^{-1}(q)=m_{A}\) ，从而 q ⊃ m \(_{A}\) ，最后 q = m \(_{A}\) B。于是环 B 是局部的。由此立即得到引理 1。

定义 1. --- 设 A 是一个局部环。若 A-代数 B 同构于 A-代数 A{]}X{[}，或者存在 A{[}X{]} 中的一个首一多项式 P，其在 \(\kappa_{\mathrm{A}}[\mathrm{X}]\) 中的像不可约，使得 B 同构于 A-代数 A{[}X{]}/(P)，则称 A-代数 B 是 A 的一个初等膨胀。

设 B 是 A 的一个初等膨胀。由前面的讨论可得下列性质：

\begin{enumerate}
\def\labelenumi{\alph{enumi})}
\item
  环 B 是局部的，且从 A 到 B 的典范同态是局部且平坦的，特别还是单射的（I，§ 3，第 5 目，命题 8）。
\item
  B 的剩余域 \(\kappa_{B}\) 是 A 的剩余域 \(\kappa_{A}\) 的一个单生成扩张。若 \(\kappa_{B}\) 在 \(\kappa_{A}\) 上的次数 d 有限，则 B 是秩为 d 的自由 A-模。
\item
  我们有 \(m_{B} = m_{A}B\) 。特别地，若 A 是域，则 B 也是域。一个域扩张是初等膨胀当且仅当它是单生成的。
\item
  若 A 是 Noether 的，则 B 也是 Noether 的。
\end{enumerate}

定义 2. --- 设 A 是一个局部环。若存在一个具有最大元 ω 的良序集 Λ，以及 B 的子代数的一个递增族 (B \(_{λ}\) ) \(_{λ∈Λ}\) ，满足下列条件，则称 A-代数 B 是 A 的一个膨胀：

\begin{enumerate}
\def\labelenumi{\alph{enumi})}
\item
  我们有 \(\mathbf{B}_{\omega} = \mathbf{B}\) ，且环 \(\mathbf{B}_{\lambda}\) 对所有 \(\lambda \in \Lambda\) 都是局部的。
\item
  若 \(\alpha\) 是 \(\Lambda\) 的最小元，则 A-代数 \(\mathbf{B}_{\alpha}\) 同构于 \(\mathbf{A}\) 。
\item
  设 \(v \neq \alpha\) 在 \(\Lambda\) 中，并设 \(S_{v}\) 是由 \(\lambda \in \Lambda\) 中满足 \(\lambda < v\) 者组成的集合。若 \(S_{v}\) 没有最大元，我们有 \(B_{v} = \bigcup_{\lambda \in S_{v}} B_{\lambda}\) ；若 \(S_{v}\) 有最大元 \(\mu\) ，则 \(B_{v}\) 是 \(B_{\mu}\) 的一个初等膨胀。
\end{enumerate}

设 B 是一个环，\(\rho: A \to B\) 是一个环同态。若由 \(\rho\) 定义的 A-代数具有该性质，则称 \(\rho\) 是一个膨胀（相应地，初等膨胀）。若如此，则 \(\rho\) 是单射。

例. --- 1) 每个域扩张都是一个膨胀。事实上，设 K 是域 k 的一个扩张。给 K 赋予一个以 0 为最大元的良序，且对 \(\lambda \in \mathbf{K}\) ，设 \(\mathbf{K}_{\lambda}\) 为 K 的一个子 \(k\) -扩张，它由 K 的元素 \(\beta\) （满足 \(\beta < \lambda\) ）生成。条件 \(a)\) ，\(b)\) ，\(c)\) 对于 \(k\) ，K 以及族 \((\mathbf{K}_{\lambda})_{\lambda \in \mathbf{K}}\) 的验证是直接的。

\begin{enumerate}
\def\labelenumi{\arabic{enumi})}
\setcounter{enumi}{1}
\tightlist
\item
  设 A 是一个局部环，I 是一个指标集。用 A{]} \((X_{i})_{i\in I}\) {[} 表示多项式环 A \([(X_{i})_{i\in I}]\) 在素理想 \(m_{A}A[(X_{i})_{i\in I}]\) 处的局部环。A-代数 A{]} \((X_{i})_{i\in I}\) {[} 是 A 的一个膨胀。事实上，给集合 I 赋予良序；设 Λ 是由 I 添加一个最大元 ω 而得到的良序集。对 i ∈ I，把 A{]} \((X_{j})_{j<i}\) {[} 等同于 B \(_{i}\) ，后者是 B = A{]} \((X_{i})_{i\in I}\) 的一个子代数，并令 B \(_{\omega}\) = B。族 (B \(_{\lambda}\) ) \(_{\lambda\in\Lambda}\) 满足条件 a)，b)，c)。
\end{enumerate}

注记. --- 采用定义 2 的记号，环 \(\mathbf{B}_{\mu}\) 是 \(\mathbf{B}_{\lambda}\) 的一个膨胀，当 \(\lambda \leqslant \mu\) 时。

命题 2. --- 设 A 是一个局部环，B 是 A 的一个膨胀。

\begin{enumerate}
\def\labelenumi{\alph{enumi})}
\item
  环 B 是局部的，且我们有 \(\mathfrak{m}_{\mathrm{A}}\mathbf{B} = \mathfrak{m}_{\mathrm{B}}\) 。
\item
  A-代数 B 是忠实平坦的。
\item
  典范同态
\end{enumerate}

\[
\gamma_ {\mathrm{B}}: \operatorname{gr} (\mathrm{A}) \otimes_ {\kappa_ {\mathrm{A}}} \kappa_ {\mathrm{B}} \rightarrow \operatorname{gr} (\mathrm{B})
\]

是双射。

\begin{enumerate}
\def\labelenumi{\alph{enumi})}
\setcounter{enumi}{3}
\tightlist
\item
  若 A 是 Noether 的，则 B 也是 Noether 的，并且 A 与 B 的 Hilbert-Samuel 级数（VIII，§ 4，第 3 目）相等。
\end{enumerate}

设 \((\mathbf{B}_{\lambda})_{\lambda\in\Lambda}\) 是 B 的子代数的一个族，满足定义 2 的条件 a)，b)，c)。

设 \(\Lambda'\) 是由满足下列性质的指标 \(\lambda \in \Lambda\) 组成的集合：对满足 \(\mu \leqslant \lambda\) 且属于 \(\Lambda\) 的每个 μ，A-代数 \(B_{\mu}\) 是局部且忠实平坦的，且 \(m_{B_{\mu}} = m_A B_{\mu}\) 。假设 \(\Lambda' \neq \Lambda\) ，并设 v 是 \(\Lambda - \Lambda'\) 的最小元。我们有 \(\alpha \in \Lambda'\) ，从而 v \(\neq \alpha\) 。\(S_{v}\) 包含于 \(\Lambda'\) 。若 \(S_{v}\) 没有最大元，我们有 \(B_{v} = \bigcup_{\lambda \in S_{v}} B_{\lambda}\) ，且 v 属于 \(\Lambda'\)

  这由第 1 目的命题 1 得出。若 \(S_v\) 有最大元 \(\mu\) ，则有 \(\mu \in \Lambda'\) 且 \(B_v\) 是 \(B_\mu\) 的一个初等膨胀：由定义 1 后面的注记，我们仍有 \(v \in \Lambda'\) ，由此得到矛盾。

当 A 是 Noether 环时，可以用类似的方式证明，集合 \(\Lambda''\) （由指标 \(\lambda \in \Lambda\) 中使得环 \(\mathbf{B}_{\lambda}\) 是 Noether 者组成）等于 \(\Lambda\) 。

因此我们有 \(\omega\in\Lambda'\) ，由此得到断言 \(a\) 和 \(b\) )。当 A 是 Noether 环时，我们有 \(\omega\in\Lambda''\) ，因而 B = B \(_{\omega}\) 是 Noether 的。

断言 c) 由 a)，b) 以及 III，§ 5，第 2 目的定理 1 得出。设 A（从而 B）是 Noether 的；由于我们有

\[
\left[ \mathfrak {m} _ {\mathrm{B}} ^ {n} / \mathfrak {m} _ {\mathrm{B}} ^ {n + 1}: \kappa_ {\mathrm{B}} \right] = \left[ \mathfrak {m} _ {\mathrm{A}} ^ {n} / \mathfrak {m} _ {\mathrm{A}} ^ {n + 1}: \kappa_ {\mathrm{A}} \right]
\]

对每个 \(n \in \mathbf{N}\) 成立，故 A 与 B 的 Hilbert-Samuel 级数相等。

推论. --- 假设 A 是 Noether 的。

\begin{enumerate}
\def\labelenumi{\alph{enumi})}
\item
  我们有 \(\dim (\mathbf{A}) = \dim (\mathbf{B})\)
\item
  假设 A 是正则的，并设 \((x_{1},\ldots ,x_{n})\) 是 A 的一个坐标系。则 B 是正则的，且序列 \((x_{1}1_{\mathrm{B}},\dots,x_{n}1_{\mathrm{B}})\) 是 B 的一个坐标系。
\end{enumerate}

这由 VIII，§ 5，第 1 目的命题 1 得出。

命题 3. --- 设 A，B，C 是三个局部环，\(u:A\to B\) ，\(v:B\to C\) 是两个膨胀。则 \(v\circ u\) 是一个膨胀。

设 \((\mathrm{B}_{\lambda})_{\lambda\in\Lambda}\) 与 \((\mathrm{C}_{\mu})_{\mu\in\mathbf{M}}\) 分别是 B 的子 A-代数族和 C 的子 B-代数族，具有定义 2 的性质 a)，b)，c)。在 \(\Lambda\) 与 M 的和 N 上，考虑这样的序关系：它在 \(\Lambda\) 和 M 上诱导给定的序，并且关系 \(\lambda<\mu\) 对 \(\lambda\in\Lambda\) ，\(\mu\in M\) 成立。这是一个良序关系。对 \(\lambda\in\Lambda\subset N\) ，令 \(C_{\lambda}=v(\mathrm{B}_{\lambda})\) 。则族 \((\mathrm{C}_{\nu})_{\nu\in\mathbf{N}}\) 满足关于 A-代数 C 的定义 1 的条件 a)，b)，c)。

定理 1. --- 设 \(f:A\to A'\) 是局部环之间的一个满射局部同态，B' 是 A' 的一个膨胀。则存在 A 的一个膨胀 B 以及从 B \(\otimes_{A}A'\) 到 B' 上的一个 A-代数同构。

\begin{enumerate}
\def\labelenumi{\Alph{enumi})}
\tightlist
\item
  假设 B' 是 A' 的一个初等膨胀。我们区分两种情形：1) 若 B' 在 A' 上有限，选取一个 A'-代数同构 \(\varphi: A'[X]/(P') \to B'\) ，其中 P' ∈ A'{[}X{]} 是一个首一多项式，其像在 \(\kappa_{A'}[X]\) 中不可约。选取一个首一多项式 P ∈ A{[}X{]} ，它在 A'{[}X{]} 中的像为 P'。它对 A 的极大理想取模必然不可约。令 B = A{[}X{]}/(P)。A-代数 B 是 A 的一个初等膨胀，且 \(\varphi\) 诱导从 \(B \otimes_{A} A'\) 到 B' 上的一个 A-代数同构。
\end{enumerate}

\begin{enumerate}
\def\labelenumi{\arabic{enumi})}
\setcounter{enumi}{1}
\tightlist
\item
  若 B' 在 A' 上不有限，选取一个 A'-代数同构 \(\psi : A']X[\to B'\) 。令 \(B = A]X[\) 。A-代数 B 是 A 的一个初等膨胀，且 B \(\otimes_{A} A'\) 典范地同构于 A'{]}X{[}。于是 \(\psi\) 诱导从 B \(\otimes_{A} A'\) 到 B' 上的一个 A-代数同构。
\end{enumerate}

\begin{enumerate}
\def\labelenumi{\Alph{enumi})}
\setcounter{enumi}{1}
\tightlist
\item
  现在过渡到一般情形。设 \((\mathbf{B}_{\lambda}^{\prime})_{\lambda\in\Lambda}\) 是由 \(A^{\prime}\) -代数组成的一个族，这些代数是 \(B^{\prime}\) 的子代数，它相对于 \(A^{\prime}\) 和 \(B^{\prime}\) 具有定义 2 的性质 a)，b)，c)。我们将用超限归纳定义一个归纳系统 \((\widetilde{\mathbf{B}}_{\lambda}, i_{\mu\lambda})\) ，它关于 \(\Lambda\) ，由局部环和单射的局部同态组成，以及同构 \(u_{\lambda}: \widetilde{B}_{\lambda} \otimes_{A} A^{\prime} \to B_{\lambda}^{\prime}\) ，使得对 \(\lambda \leqslant \mu\) ，\(u_{\mu} \circ (i_{\mu\lambda} \otimes \operatorname{Id}_{A^{\prime}}) \circ u_{\lambda}^{-1}\) 是 \(B_{\lambda}^{\prime}\) 到 \(B_{\mu}^{\prime}\) 中的典范单射。
\end{enumerate}

若 \(\alpha\) 是 \(\Lambda\) 的最小元，我们令 \(\tilde{\mathbf{B}}_{\alpha} = \mathbf{A}\) ，\(i_{\alpha \alpha} = \mathrm{Id}_{\mathbf{A}}\) ，并取 \(u_{\alpha}\) 为典范同构 \(\mathbf{A} \otimes_{\mathbf{A}} \mathbf{A}' \to \mathbf{A}'\) 。

设 \(v \in \Lambda\) ，并假设 \(\tilde{\mathbf{B}}_{\lambda}\) ，\(u_{\lambda}\) 和 \(i_{\mu \lambda}\) 当 \(\lambda \leqslant \mu < v\) 时均已构造。设 \(S_v\) 是由元素 \(\varepsilon\) 组成的集合，这些元素属于 \(\Lambda\) 且满足 \(\varepsilon < v\) 。若 \(S_v\) 没有最大元，则取 \(\tilde{\mathbf{B}}_v\) 为诸 \(\tilde{\mathbf{B}}_{\lambda}\) （\(\lambda \in S_v\) ）的归纳极限，取 \(u_v\) 为合成同构 \(\tilde{\mathbf{B}}_v \otimes_A A' \to \varinjlim (\tilde{\mathbf{B}}_{\lambda} \otimes_A A') \to \varinjlim B'_\lambda \to B'_v\) ，而取 \(i_{v\lambda}\) （当 \(\lambda \in S_v\) 时）为从 \(\tilde{\mathbf{B}}_{\lambda}\) 到 \(\tilde{\mathbf{B}}_v\) 的典范映射。若 \(S_v\) 有最大元 \(\mu\) ，则 \(B'_v\) 是 \(B'_\mu\) 的一个初等膨胀。由 A)，存在一个初等膨胀 \(i_{\nu \mu}:\tilde{\mathbf{B}}_{\mu}\to \tilde{\mathbf{B}}_{\nu}\) 以及一个 \(\tilde{\mathbf{B}}_{\mu}\) -代数同构，它把 \(\tilde{\mathbf{B}}_{\nu}\otimes_{\tilde{\mathbf{B}}_{\mu}}\mathbf{B}_{\mu}^{\prime}\) 映到 \(\mathbf{B}_{\nu}^{\prime}\) 上。取 \(u_{\nu}\) 为合成而得的 A-代数同构

\[
\widetilde {\mathbf {B}} _ {\nu} \otimes_ {\mathrm{A}} \mathrm{A} ^ {\prime} \rightarrow \widetilde {\mathbf {B}} _ {\nu} \otimes_ {\widetilde {\mathbf {B}} _ {\mu}} (\widetilde {\mathbf {B}} _ {\mu} \otimes_ {\mathrm{A}} \mathrm{A} ^ {\prime}) \rightarrow \widetilde {\mathbf {B}} _ {\nu} \otimes_ {\widetilde {\mathbf {B}} _ {\mu}} \mathrm{B} _ {\mu} ^ {\prime} \rightarrow \mathrm{B} _ {\nu} ^ {\prime}
\]

而取 \(i_{\nu \lambda}\) （当 \(\lambda \in S_{v}\) 时）为同态 \(i_{\nu \mu} \circ i_{\mu \lambda}\) 。

然后令 \(\mathbf{B} = \tilde{\mathbf{B}}_{\omega}\) ，且对所有 \(\lambda \in \Lambda\) ，用 \(\mathbf{B}_{\lambda}\) 表示 \(\tilde{\mathbf{B}}_{\lambda}\) 在典范单射 \(\tilde{\mathbf{B}}_{\lambda} \to \mathbf{B}\) 下的像。族 \((\mathbf{B}_{\lambda})_{\lambda \in \Lambda}\) 满足定义 2 的条件 \(a)\) ，\(b)\) ，\(c)\) ，且 \(\mathbf{B}\) 是 A 的一个膨胀。另一方面，同态 \(u_{\omega}\) 是从 \(\mathbf{B} \otimes_{A} \mathbf{A}'\) 到 \(\mathbf{B}'\) 上的一个 A'-同构。

推论. --- 设 A 是一个局部环，K 是其剩余域 \(\kappa_{\mathrm{A}}\) 的一个扩张。则存在一个局部环 B 以及一个膨胀 \(\mathbf{A} \to \mathbf{B}\) ，使得 \(\kappa_{\mathrm{A}}\) -代数 \(\kappa_{\mathrm{B}}\) 同构于 K。

事实上，同态 \(\kappa_{\mathbf{A}} \to \mathbf{K}\) 是一个膨胀（例 1）。取 \(\mathbf{A}' = \kappa_{\mathbf{A}}\) 和 \(\mathbf{B}' = \mathbf{K}\) 应用定理 1，我们得到 A 的一个膨胀 B 的存在性，以及从 \(\mathbf{B} / \mathfrak{m}_{\mathbf{A}}\mathbf{B}\) 到 K 上的一个 A-同构，由此得到本推论。

\refstepcounter{appendixitem}
\subsection*{\theappendixitem.\ p-环的存在性}
\phantomsection
\addcontentsline{toc}{subsection}{\theappendixitem.\ p-环的存在性}

命题 4. --- 设 \(p\) 是一个素数，\(k\) 是一个特征为 \(p\) 的域，并设 \(n\) 是一个整数 \(\geqslant 1\) ，或为 \(+\infty\) 。则存在一个 \(p\) -环（§ 2，第 1 目，定义 1），其长度为 \(n\) ，剩余域同构于 \(k\) 。

可以把 \(k\) 看作 \(\mathbf{Z}/p\mathbf{Z}\) （局部环 \(\mathbf{Z}_{(p)}\) 的剩余域）的一个扩张。由定理 1 的推论，存在一个局部环 \(\mathbf{B}\) ，它是 \(\mathbf{Z}_{(p)}\) 的一个膨胀，使得 \(\kappa_{\mathrm{B}}\) 同构于 \(k\) 。局部环 \(\mathbf{Z}_{(p)}\) 是正则的，且 \(\{p\}\) 是 \(\mathbf{Z}_{(p)}\) 的一个坐标系。由第 2 目命题 2 的推论，环 \(\mathbf{B}\) 是正则的，且 \(\{p1_{\mathrm{B}}\}\) 是 \(\mathbf{B}\) 的一个坐标系。换言之，\(\mathbf{B}\) 是一个离散赋值环，其极大理想为 \(p\mathbf{B}\) 。于是完备化 \(\mathbf{C}\) （由 \(\mathbf{B}\) 完备化而得）是一个无限长度的 \(p\) -环，且剩余域 \(\kappa_{\mathrm{C}}\) 同构于 \(\kappa_{\mathrm{B}}\) ，从而同构于 \(k\) 。此外，对每个整数 \(n \geqslant 1\) ，C/ \(p^{n}\) C 是一个 \(p\) -环，其长度为 \(n\) ，剩余域同构于 \(\kappa_{\mathrm{C}}\) ，从而同构于 \(k\) 。

\specialchapter{习题}
\setcounter{exercisesection}{0}
